% template_explorer_guide.tex -- a starting point for a NEW student guide in the SCA101 house style.
%
% Copy this file, rename it, replace the text, and build:   pdflatex my_guide.tex   (twice, so the contents page fills in)
% Every macro and environment of scaguide.sty is used once below, with a comment saying what it is for.
% The three finished guides next to this file (Lab_2_1A_Explorer_Guide.tex, Lab_2_1B_Explorer_Guide.tex,
% Lab_2_1A_Real_Husky_Guide.tex) are real examples of the same pieces in use.
\documentclass[11pt,letterpaper]{article}
\usepackage{scaguide}

% header (top left / top right), footer (bottom left) and the PDF's title and subject
\guidesetup{
  headerleft={ChipWhisperer SCA101  -  Lab X Explorer Guide},
  headerright={My New Explorer},
  footerleft={Companion to: Lab X - My New Explorer.ipynb},
  title={SCA101 Lab X Explorer - Step-by-Step Student Guide},
  subject={Step-by-step guide to the Lab X Explorer notebook}}

\begin{document}

% ---------------------------------------------------------------- cover: kicker, title (use \\ for a line break), subtitle
\guidecover{SCA101  -  LAB X  (EXPLORER)}{My New\\Explorer}{Step-by-step student guide: edit the numbers, watch the result change}

Write one or two paragraphs here that say what the lab is about and what the notebook lets the student change. Plain text, \textbf{bold} for
the key idea, \emph{italics} for a term, and \code{inline code} for anything the student types or sees on the screen
(write \verb|\_| for an underscore inside normal text, for example \code{MY\_VARIABLE}).

% ---------------------------------------------------------------- call-out boxes: kind = do | see | ifnot | warn | tip
\begin{guidecallout}{tip}{HOW TO USE THIS GUIDE}
This guide goes \textbf{with} the notebook. Keep the notebook open next to this PDF.
\end{guidecallout}

\guidehthree{The four colored boxes you will see everywhere}   % small teal heading (not in the contents)
\guidelegend                                                   % the four-box legend

\clearpage
\tableofcontents    % filled in by the \guidehone / \guidehtwo / \guidestep headings; run pdflatex twice
\clearpage

% ---------------------------------------------------------------- headings
%   \guidehone{...}  chapter heading        \guidehtwo{...}  section heading
%   \guidestep{N}{...}  numbered step: navy badge + heading      \guidehthree{...}  small teal heading
\guidehone{The big picture}
What is the attack, or the idea? What is different about the Explorer?

\guidehone{Step 1 - look at the data}
\guidestep{1}{Look at a trace}
Explain what the cell does. Then list the variables the student may edit:

% ---------------------------------------------------------------- tables
% Columns are R{fraction of the text width}; the fractions should add up to 1. The header row is navy; rows alternate white/grey.
% Use \codesmall{...} for code inside a table cell and \textbf{...} for bold.
{\small\arrayrulecolor{guidegrid}
\rowcolors{2}{white}{guidezebra}
\begin{longtable}{|R{0.22}|R{0.18}|R{0.60}|}
\hline
\rowcolor{guidenavy}\hdrcell{Variable} & \hdrcell{Default} & \hdrcell{What it does} \\ \hline
\endhead
\codesmall{TRACE}      & \codesmall{2}     & Which recorded trace to show: 0, 1, 2 or 3. \\ \hline
\codesmall{START, END} & \codesmall{0, 600} & The zoom window, in samples. \\ \hline
% a table cell with a source tag (used by the troubleshooting tables):  text\srctag{Reproduced}
\end{longtable}}

% ---------------------------------------------------------------- figures: \guidefigure{file}{width as a fraction of the text width}{caption}
\guidefigure{figures/hw_step3.png}{0.85}{Caption: what the picture shows. Put your PNG files in the figures folder.}

% ---------------------------------------------------------------- the five kinds of call-out
\begin{guidecallout}{do}{DO THIS}
\textbf{Try:} (1) change \code{TRACE} to 3. (2) Raise \code{SMOOTH}. What changes?
\end{guidecallout}

\begin{guidecallout}{see}{YOU SHOULD SEE}
The plot redraws and the line under it says \code{showing samples 0 to 600}.
\end{guidecallout}

\begin{guidecallout}{ifnot}{IF IT DOESN'T}
\code{FileNotFoundError: ...} -- the traces folder is missing. Run the Setup cell again.
\end{guidecallout}

\begin{guidecallout}{warn}{WATCH OUT}
A mistake that costs time, and how to avoid it.
\end{guidecallout}

\begin{guidecallout}{tip}{GOOD TO KNOW}      % the title is free text: use it to name the box
Extra background that is nice to have but not needed to finish the step.
\end{guidecallout}

% ---------------------------------------------------------------- code blocks (verbatim, keep lines under about 100 characters)
\begin{guidecode}
python -m pip install cwtraces "numpy<=1.26.4"
\end{guidecode}

% ---------------------------------------------------------------- lists
\begin{enumerate}[label=\arabic*.,leftmargin=1.8em,itemsep=2pt,topsep=2pt,parsep=0pt]
\item Start Jupyter and open the notebook.
\item Choose \textbf{Kernel > Restart \& Run All}.     % write \& for an ampersand
\end{enumerate}

\guidehtwo{You are done when...}
\begin{itemize}[leftmargin=1.6em,itemsep=2pt,topsep=2pt,parsep=0pt]
\item You can say what the trace shows and why.
\end{itemize}

% keep a block together on a page / start a new page if less than this much room is left
\Needspace{120pt}

% ---------------------------------------------------------------- the licence notice that every guide ends with
\guidehthree{Notice}
\guidesmall{This guide accompanies ChipWhisperer course material that references (commercial) training material on ChipWhisperer.io. You can freely execute
and use the lab per the open-source license (including in your own courses if you distribute similarly), but you must maintain notice about the source:
https://www.ChipWhisperer.io.}

\end{document}
